% 06-aplicaciones.tex — Chapter 6: Aplicaciones Financieras

\chapter{Aplicaciones Financieras}
\label{ch:aplicaciones}

Aunque el modelo SF-Harris se reduce a sus ingredientes esenciales, la
distribuci\'{o}n emp\'{\'i}rica con estimaci\'{o}n bayesiana tiene aplicaciones
directas en finanzas.

\section{Pricing de swaps de varianza}

Un swap de varianza paga $\text{RV}_T - K_{\text{var}}$ al vencimiento, donde
$K_{\text{var}}$ es el \emph{strike} de varianza. El strike justo es:
\begin{equation}
    K_{\text{var}} = \E[\RV_T] = \E[\tau_T^*]
\end{equation}

La distribuci\'{o}n emp\'{\'i}rica $Q_{\text{emp}}$ estima directamente
$\E[\tau^*]$ como la media de $\RV_t$ en la ventana de entrenamiento.
La ventaja sobre modelos param\'{e}tricos es que $Q_{\text{emp}}$ captura
la asimetr\'{\'i}a y curtosis sin imponer una forma funcional.

Para pasar de la medida f\'{\'i}sica $\Prob$ a la neutral al riesgo $\Q$,
se a\~{n}ade una prima de riesgo de volatilidad $\lambda_v$:
\begin{equation}
    \E^\Q[\RV_T] = e^{\lambda_v T} \cdot \E^\Prob[\RV_T]
\end{equation}
donde $\lambda_v$ se calibra con opciones liquidas~\citep{demeterfi1999}.

\section{Opciones sobre VIX}

El VIX es la volatilidad impl\'{\'i}cita a 30 d\'{\'i}as extra\'{\'i}da de opciones
sobre el S\&P~500. Nuestra distribuci\'{o}n predictiva de $\tau^*$ proporciona
intervalos de probabilidad calibrados para $\text{VIX}_T$:

\begin{equation}
    P(\text{VIX}_T \leq v) \approx P\left(\sqrt{\frac{252}{T} \sum_{t=1}^T \tau_t^*} \leq v\right)
\end{equation}

La distribuci\'{o}n emp\'{\'i}rica captura la asimetr\'{\'i}a positiva del VIX
(sin VIX negativo, cola derecha pesada) sin suponer normalidad.

\section{Trading del skew de volatilidad}

El \emph{skew} de volatilidad es la diferencia entre volatilidad impl\'{\'i}tica
de puts out-of-the-money y calls out-of-the-money. La distribuci\'{o}n
emp\'{\'i}rica $Q_{\text{emp}}$ tiene asimetr\'{\'i}a negativa (colas izquierdas
m\'{a}s pesadas para $\log(\RV)$), lo que genera un skew natural.

La se\~{n}al de trading consiste en comparar el skew impl\'{\'i}cito del mercado
con el skew predicho por $Q_{\text{emp}}$:
\begin{itemize}
    \item Si el skew impl\'{\'i}cito es m\'{a}s plano que el de $Q_{\text{emp}}$,
          vender straddles y comprar strangles.
    \item Si el skew impl\'{\'i}cito es m\'{a}s pronunciado, la operaci\'{o}n inversa.
\end{itemize}

\section{Estructura temporal de $\pstay$}

Aunque $\pstay$ no es informativo para la cobertura de $\log(\RV)$, la
probabilidad de estancia \textbf{var\'{\'i}a con la frecuencia de muestreo}:
\begin{itemize}
    \item 15\,min: $\pstay \approx 0.88$ (estimado por ACF)
    \item Diario: $\pstay \approx 0.0$ (pr\'{a}cticamente iid)
\end{itemize}

Esto tiene una interpretaci\'{o}n financiera: a frecuencia intradiaria, el
precio ``se queda'' en niveles de volatilidad similares entre intervalos cortos.
A escala diaria, los cambios son iid respecto a la distribuci\'{o}n marginal.

Para \emph{term structure modeling}, la curva $\pstay(f)$ podr\'{\'i}a calibrar
la correlaci\'{o}n entre volatilidades a diferentes horizontes, pero nuestra
evidencia muestra que esta informaci\'{o}n no mejora la cobertura.

\section{Trading de dispersi\'{o}n}

La dispersi\'{o}n (\emph{dispersion trading}) explota la diferencia entre la
volatilidad impl\'{\'i}cita del \'{\'i}ndice y la volatilidad impl\'{\'i}cita
de sus componentes. La distribuci\'{o}n emp\'{\'i}rica $Q_{\text{emp}}$ aplicada a
cada componente estima la covarianza de la matriz de varianza-covarianza:
\begin{equation}
    \widehat{\text{Cov}}(\RV_i, \RV_j) = \frac{1}{n} \sum_{t=1}^{n} (\RV_{i,t} - \bar{\RV}_i)(\RV_{j,t} - \bar{\RV}_j)
\end{equation}

La covarianza emp\'{\'i}rica es consistente con nuestro enfoque: la distribuci\'{o}n
conjunta emp\'{\'i}rica captura correlaciones sin suponer normalidad multivariada.

\section{Lo que NO funciona}

Es igualmente importante se\~{n}alar las aplicaciones donde el modelo no aporta valor:

\begin{enumerate}
    \item \textbf{Predicci\'{o}n direccional}: El modelo predice la
          \emph{distribuci\'{o}n} de la volatilidad, no su direcci\'{o}n.
          Los rendimientos condicionales $r_t \mid \tau_t \sim N(0, \tau_t)$
          tienen media cero por construcci\'{o}n.

    \item \textbf{Detecci\'{o}n de crisis}: La sorpresa predictiva tiene
          AUC = 0.416 (anti-predictivo). El modelo siempre espera colas pesadas,
          as\'{\'i} que las crisis no son sorprendentes bajo el modelo.

    \item \textbf{Predicci\'{o}n diaria}: A escala diaria, el ruido en $\hat{\tau}_t$
          domina (AAD 2--12\,pp). Se requieren datos intradiarios para que el
          modelo funcione.
\end{enumerate}